\documentclass[10pt,letterpaper]{article}
\usepackage{cornell-notes}

\title{Chapter 18: Spanning Forest}
\author{}
\date{}

\setDocTitle{Chapter 18: Spanning Forest}
\setDocSubtitle{Combinatorial Algorithms}
\setDocVersion{v1.0}
\setDocStatus{Study Companion}
\setDocOwner{Combinatorial Algorithms Cornell Notes}
\setDocAuthor{}
\setDocDate{\today}
\setDocSummary{Cornell-format study and review notes for Chapter 18: Spanning Forest.}

\setCornellCollection{Combinatorial Algorithms Cornell Notes}
\setCornellUnitType{Chapter}
\setCornellUnitNumber{18}
\setCornellUnitTitle{Spanning Forest}
\setCornellCompanionLabel{Study and review companion}

\hypersetup{pdftitle={Chapter 18: Spanning Forest},pdfsubject={Cornell notes},pdfauthor={}}

\begin{document}
\maketitle
\begin{tcolorbox}[colback=Paper,colframe=Ink]{\small\bfseries\color{Teal}CORNELL NOTES}\par{\LARGE\bfseries Chapter 18: Spanning Forest of a Graph (SPANFO)}\par\medskip\textbf{Topic:} Connectivity, traversal, and compact edge rearrangement\hfill\textbf{Date:} \today\par\textbf{Study purpose:} Produce component labels and a rooted spanning-forest prefix.\end{tcolorbox}
\begin{tcolorbox}[colback=Teal!5,colframe=Teal,title=Essential question]How can an edge-list graph be rearranged so its spanning forest comes first while component labels are produced in linear edge work?\end{tcolorbox}
\begin{tcolorbox}[colback=white,colframe=Mist,title=Learning objectives]\begin{itemize}\item State all four output requirements.\item Contrast stack-based depth-first and queue-like breadth-first approaches.\item Explain linked incidence lists encoded inside the edge array.\item Trace LINK, VISIT, SCAN, COMPONENT, and SPANFO.\item Recognize the destructive-storage contract.\end{itemize}\textbf{Prerequisites:} undirected graphs, components, spanning trees, DFS/BFS, linked lists.\end{tcolorbox}
\section*{Cornell notes}
\begin{longtable}{@{}Q!{\color{Mist}\vrule width 1pt}A@{}}\cornellhead
\cue{What must be returned?}{The component count $k$; a component label $x(i)$ for every vertex; the edge list rearranged with all spanning-tree edges first and nonforest edges last; and an orientation/order in which every forest edge points away from a root and root-to-vertex walks use increasing edge indices.}
\cue{What does depth-first search show?}{A stack stores newly discovered vertices. First encounters assign component and scan numbers and produce tree edges. For an already assigned neighbor, predecessor and scan-order tests distinguish the parent edge, a first encounter with an extra edge, and the duplicate encounter from the opposite endpoint.}
\cue{Why implement a breadth-first variant?}{The selected routine reduces named array storage to the input edge array and one vertex array. Vertices are scanned in visit order, behaving like a queue, while the edge array temporarily carries linked incidence lists and final tree-address information.}
\cue{What does LINK build?}{Replace each occurrence of vertex $p$ in the endpoint array by a pointer to the previous occurrence, and let $x(p)$ point to the latest occurrence. Encode endpoint addresses with a base $M=1+\max(n,E)$ and use signs to distinguish unvisited, visited, scanned, and restored states.}
\cue{What does VISIT decide?}{Starting from one endpoint occurrence, follow the neighbor's linked list. It either discovers an unvisited neighbor and marks the list, abandons because that vertex was already visited, or reports an edge to a previously scanned neighbor that can attach the current vertex to the growing tree.}
\cue{How do SCAN and COMPONENT cooperate?}{SCAN walks all incidence entries of a vertex, restores endpoint names, links newly discovered vertices into scan order, and retains one suitable tree edge. COMPONENT repeatedly scans the next queued vertex, increments the component count for a root, and records each selected tree edge with a future forest position.}
\cue{What does the final SPANFO pass do?}{After all components are traversed, negative address marks identify selected tree edges. Swap each selected edge into its assigned prefix location, orient it from the nearer-root endpoint to the farther endpoint, and translate the temporary vertex information into final component numbers.}
\cue{Correctness and cost}{Every vertex is visited once, every incidence is followed a bounded number of times, and exactly one tree edge attaches each nonroot vertex. Thus each component contributes $|V_c|-1$ tree edges and contains no cycle. The source describes linear-time offline traversal in $E$ with storage encoded in $2E+n$ words.}
\cue{Cautions}{The endpoint array is destroyed and rearranged. Pointer packing assumes a signed integer range large enough for multiples of $M$. Loops, parallel edges, isolated vertices, and disconnected graphs require the exact visit tests. A modern implementation may prefer explicit adjacency and queue structures unless destructive compact storage is required.}
\cue{Retrieval practice}{\begin{enumerate}\item What four outputs are required?\item Why are scan numbers useful in DFS?\item What does LINK encode?\item How is one tree edge selected for a nonroot?\item Why does the forest prefix contain $n-k$ edges?\end{enumerate}}
\end{longtable}
\begin{tcolorbox}[colback=Ink!4,colframe=Ink,title=Cornell summary]
SPANFO is an offline connectivity routine with a strong output contract: component labels, component count, a spanning-forest prefix, and root-consistent edge order. The exposition first gives a depth-first solution, then implements a breadth-first variant that stores linked incidence lists inside the mutable endpoint array. LINK builds those lists; VISIT classifies neighbors; SCAN processes a vertex; COMPONENT processes one connected component; and a final pass moves marked tree edges into their assigned prefix positions. Each nonroot vertex receives exactly one attachment edge, establishing a cycle-free spanning tree per component. The compact representation gives linear edge work but deliberately destroys the original edge order and depends on carefully packed signed pointers.
\end{tcolorbox}
\end{document}
